\origpage{28}{}
For $i = 1, 2$ let $\Delta_i \subseteq \Sigma$ be the closed disc bounded by $F \cup J_i$. Now $(\Gamma \cap \Delta_j) \cup \mu(J_{\delm{i}\fix{j}})$\fixnote{fix:6.7-proof}{Here $\Delta_j$ is bounded by $F \cup J_j$, so the boundary graph that goes with it is $\mu(J_j)$. The index $i$ is not bound here, and for $i \ne j$ the graph $\mu(J_i)$ lies on the other disc.} has a presentation in $\Sigma$ with depth $\le q$ and defect $< r$, and it has an exhaustive vital $(p + 2q)$-linkage, and so $\Delta_j$ includes at most $\lambda(p + 2q, q)$ $\Sigma$-jumps, from the inductive hypothesis. But every $\Sigma$-jump belongs to one of $\Delta_1, \Delta_2$, and so there are at most $2\lambda(p + 2q, q) = \lambda$ $\Sigma$-jumps altogether, as required. $\square$

\textbf{Proof of (6.3).} We proceed by induction on $\Sigma$ and by \del{(6.3)}\fix{(6.7)}\fixnote{fix:6.3-proof}{The reference meant is (6.7): a proof of (6.3) cannot use (6.3), and (6.7) is the case of a disc.}, we may assume that $\Sigma$ is not a disc. We make the inductive hypothesis that

\hspace*{12pt}(1) \textit{For all integers $p, r \ge 0$ and $q \ge 1$ there exists $\lambda(p, q, r) \ge 0$ such that for every surface $\Sigma'$, if either $\hat{\Sigma}'$ is simpler then $\hat{\Sigma}$, or $\hat{\Sigma}'$ is homeomorphic to $\hat{\Sigma}$ and $c(\Sigma') < c(\Sigma)$, the following is true. If $G'$ is a graph with a presentation in $\Sigma'$ of depth $\le q$ and defect $\le r$, and there is an exhaustive vital $p$-linkage in $G'$, then there are at most $\lambda(p, q, r)$ $\Sigma'$-jumps.}

Let $\lambda = \max(\lambda(p + 4q, q, r + 2), 2\lambda(p + \delm{2q}\fix{4q}, q, r + 1))$. We claim that $\lambda$ satisfies the theorem. For let $\Gamma, F_1^*, \dots, F_r^*, H_1, \dots, H_r, \mu$ be a presentation of a graph $G$ in $\Sigma$, with depth $\le q$.

\hspace*{12pt}(2) \textit{We may assume that every proper line $F \subseteq \Sigma$ with ends in $V(\Gamma)$ and with no internal point in $U(\Gamma)$ separates $\Sigma$.}

\textit{Subproof.} Suppose that $F \subseteq \Sigma$ is a proper line with ends $a, b \in V(\Gamma)$ and with no internal point in $U(\Gamma)$, and $F$ does not separate $\Sigma$. Let $\Sigma'$ be obtained from $\Sigma$ by cutting along $F$. By (4.1)(i) and (4.1)(ii), either $\hat{\Sigma}'$ is simpler than $\hat{\Sigma}$, or $\hat{\Sigma}'$ is homeomorphic to $\hat{\Sigma}$ and $c(\Sigma') < c(\Sigma)$. Let $\Gamma'$ be the drawing in $\Sigma'$ obtained from $\Gamma$ by splitting appropriately the vertices of $\Gamma$ in $F$. Let $a \in F_1^*, b \in F_2^*$ say\fix{, or let $a, b \in F_1^*$; in the second case the splitting below is done twice inside $F_1^*$, first at $a$ and then at $b$}.\fixnote{fix:6.3-claim2-ends}{The ends $a, b$ of $F$ may lie in the same support, a case that ``say'' does not cover. For instance, a non-separating line in an annulus may start and end on one cuff that is a single support. Splitting twice inside the support gives each vertex of $V(\mu(a)) \cap V(\mu(b))$ three copies. Each split is one use of (2.4) per vertex, so the linkage is still a $(p + 4q)$-linkage, and the defect is still $\le r + 2$. So the bound of Claim (2) stands.} Let $F_1^*$ have ends $u_1, u_2$, and let $I_i \subseteq F_1^*$ be the support with ends $u_{\delm{1}\fix{i}}, a$.\fixnote{fix:6.3-claim2-supports}{The two supports $I_1, I_2$ have ends $u_1, a$ and $u_2, a$. With $u_1$ for both, $I_1 = I_2$, and the splitting by $\delta_1(v), \delta_2(v)$ below would split nothing.} For each $v \in V(\mu(a))$, let $\delta_i(v)$ be the set of edges of $\mu(I_i)$ incident with $v$ $(i = 1, 2)$. By splitting $v$ according to $\delta_1(v), \delta_2(v)$, for each $v \in V(\mu(a))$, and similarly splitting each $v \in V(\mu(b))$, we obtain a graph $G'$ which has a presentation in $\Sigma'$ (using the drawing $\Gamma'$) of depth $\le q$ and defect $\le r + 2$. But by (2.4), it follows that $G'$ has an exhaustive vital $(p + 4q)$-linkage, and so by (1), $\Gamma'$ has at most $\lambda(p + 4q, q, r + 2)$ $\Sigma'$-jumps. But these $\Sigma'$-jumps are in 1-1 correspondence with the $\Sigma$-jumps of $\Gamma$, and so $\Gamma$ has at most $\lambda(p + 4q, q, r + 2) \le \lambda$ $\Sigma$-jumps, as required. This proves (2).

\hspace*{12pt}(3) $c(\Sigma) = 1$.
